\scanpage{54}
% @01, underlined heading
\uline{Intégration} \uline{des Algèbres $1/2$ simpliciales}

% @02–06, graded object
Soit $A_*$ un groupe gradué différentiel $1/2$ simplicial i.e. un objet $1/2$ simplicial dans la catégorie des groupes gradués différentiels (: op. de dérivation de degré $+1$) \deleted{On suppose} Donc $\forall n\geq0$ $A_n=A_n^*=\sum_{p\in\mathbb{Z}}A_n^p$, avec op. différentiel $d_n\colon A_n^{\unsure{\bullet|*|i}}\to A_n^{\unsure{\bullet|*}}$, $A_n^p\to A_n^{p+1}$. On suppose

% @07–11, hypothesis
(1) $\forall n$, on a $A_n^p=0$ si $p\notin[0,n]$, et \deleted{\unsure{$l'$|$k$}} $A_n^*$ muni de l'augmentation \inserted{$\varepsilon_n\colon$} $A_0\to A_n^*$ [\uline{NB} associée à l'unique $\sigma\colon\Delta_n\to\Delta_0$ ; $A_0$ est considéré comme groupe doué de la graduation nulle et d'op. diff trivial] est une \uline{résolution} de $A_0$.

% @12–13, boundary
Rappelons qu'on a sur le système des $A_n$ un op. différentiel (de degré $-1$)
% @14–15, two defining lines
\[
\begin{aligned}
 \partial&\colon A_n^*\longrightarrow A_{n-1}^*\\
 \partial(\omega)&=\sum_{0\leq i\leq n}\deleted{\unsure{d}}\partial_{n,i}^*(\omega)(-1)^i
\end{aligned}
\]
% @16–17, faces
(où les $\partial_{\unsure{n,i|n-i}}$ \unsure{$1\leq i\leq n$|$0\leq i\leq n$}) sont les $n+1$ applications strictmt croissantes $\Delta_{n-1}\to\Delta_n$. On sait que $\partial(\partial\omega)=0$

% @18–20, proposition
\uline{Proposition} 1 Il existe un et un seul système d'hom de groupes
% @19, maps
\[
 A_n^*\xrightarrow{k_n}A_0
\]
% @20, conditions
satisfaisant aux conditions suivantes
% @20a, left margin
% slanted
\margin{\uline{NB} le système des $k_n$ \uline{ne commute pas} en gén. aux op. $1/2$ simpliciales !}
% @21–23, conditions
\[
\begin{aligned}
 &\text{a) (normalisation)}& k_0&=\operatorname{id}_{A_0}\\
 &\text{b) (Condition de degré)}& k_{\unsure{n|m}}(A_n^p)&=0\quad\text{si }p\ne n\\
 &\text{c) (Stokes)}& k_n(d\omega)&=k_{n-1}(\partial\omega)\quad:\ \forall\omega\in A_n^\bullet.
\end{aligned}
\]
% @24–25, reduction
En vertu de b), la condition c) équivaut à la $=$ condition, avec $\omega\in A_n^{n-1}$. On montre existence et unicité \deleted{par} \inserted{\unsure{degré $\leq N$ par|de $k_n$, $n\leq N$, par|degré $\in\mathbb N$ par}} récurrence sur $N$, en notant que pour $N=0$ c'est trivial. Si $N>0$, \deleted{\unsure{si c'est|$\in A$}} il reste à démontrer ex. et un. de $k_N$ i.e. des $k_N(\overline\omega^N)$ pour $\deleted{\unsure{N}}\overline\omega^N\in A_N^N$.
% @28–30, exactness and Stokes
Or $d\overline\omega^N=0$ \deleted{\unsure{puisque}} car $A_N^{N+1}=0$, \inserted{(1)} donc $\overline\omega^N=d\omega^{N-1}$ (1). [En fait il suffit ici \deleted{\unsure{ici|i.e.|ier}} $H^N(A_N^*)=0$ pour $N\geq1$] on le \uline{définit} par c) i.e.
% @31, last equation
\[
 k_n(d\omega^{N-1})=k_{n-1}(\partial\omega)
\]
% archive sheet number: 53

